\section{Discussion}

EHA reveals failures below answer accuracy. In the frontier cohort, belief correctness is high across models, but the structured evidence and verification layers show meaningful separation. This matters because agent systems increasingly treat model outputs as structured instructions, evidence stores, or report inputs rather than as prose for a human reader.

Generated lore is not merely a content problem. It is a provenance and evidence-role problem. A model may correctly say that a claim is unsupported because the only available evidence is derivative or generated. But if the model still places those derivative documents into \code{supporting_evidence}, it risks laundering rejected material into downstream support.

Structured fields therefore matter. A prose explanation and a JSON field can disagree in practice. A human reader may reconcile them; a downstream system may not. EHA's scoring contract makes that mismatch visible.

The prompt-hygiene result also limits a common mitigation strategy. Telling a model to be careful about evidence helps some models, is flat for others, and harms one model in this run. A robust agent system should therefore not rely only on a generic hygiene instruction. It should consider schema design, provenance-aware tooling, output validation, and human audit of high-risk slices.
